
\section{Introduction}
\label{sec:introduction}

\subsection{From ``constants as inputs'' to ``constants as protocol-stable fixed points''}
\label{sec:intro_constants_as_fixed_points}

Modern theoretical physics often starts from an action principle: a variational functional selects admissible trajectories and field configurations.
Yet two structural questions remain, even when the dynamical equations are known:
\begin{itemize}
\item \textbf{Origin of effective potentials.} Why does the potential term take the form it does, as opposed to another form consistent with symmetries?
\item \textbf{Fine structure of dimensionless constants.} Why do certain constants take the observed numerical values, rather than being free inputs of the model?
\end{itemize}
In standard presentations, these elements are usually introduced as externally specified data.

This paper adopts an alternative starting point: \emph{finite resources} and \emph{auditable protocols}.
We treat scanning and readout---finite-time averaging, finite-resolution regularization, and explicit algorithmic costs---as primary objects rather than post-hoc numerical procedures applied to an ``ontic'' continuum.
The guiding hypothesis is that selection should be sought not only at an idealized infinite limit, but in the \emph{stability and cost} of the finite process that approaches that limit.

\subsection{Wish as a minimal operative target}
\label{sec:intro_wish}

Building on the classical scan-to-integral correspondence for equidistributed sequences \cite{Weyl1916,KuipersNiederreiter1974} and on the period viewpoint \cite{KontsevichZagier2001}, we elevate a key object to axiomatic status:
\emph{Wish} as \emph{protocol-stable period data}.
The term ``Wish'' is used purely as a technical label for an auditable target object:
\begin{itemize}
\item it is \textbf{not} a psychological statement;
\item it is a \textbf{structured datum} (domain, kernel, regularization);
\item it is \textbf{protocol-stable}: invariant under protocol equivalences and stably approximable under finite resources with an explicit error budget.
\end{itemize}
In this sense, ``constants'' enter not as external numbers but as \emph{stable invariants of a protocol class}.

\subsection{Audit layers and logical scope}
\label{sec:intro_audit_layers}

We maintain a strict separation between two layers:
\begin{itemize}
\item \textbf{Closed layer.} Definitions of protocols, readouts, period data, functorial interfaces, and deterministic error certificates; and theorems stated and proved solely within this vocabulary.
\item \textbf{Programmatic layer.} Interpretive statements and model-building suggestions (e.g.\ ``constant selection'' and teleological language) appear only as falsifiable conjectures or as mappings between closed objects.
\end{itemize}
No theorem in the closed layer depends on interpretive assumptions.

\paragraph{Reader's map.}
Closed-layer claims include the realization theorem (Theorem~\ref{thm:period_realization}), certified error bounds (Theorem~\ref{thm:koksma_hlawka} and \eqref{eq:certified_total_error}), explicit discrepancy certificates (equation \eqref{eq:golden_discrepancy_bound} and Theorem~\ref{thm:etk}), and the variational dynamics with Lyapunov monotonicity (Theorem~\ref{thm:teleology_dynamics} and Proposition~\ref{prop:energy_decay}, including the nonsmooth certificate form in Proposition~\ref{prop:nonsmooth_energy}).
Programmatic claims are explicitly labeled as conjectures (notably Conjecture~\ref{conj:selection}) or confined to Section~\ref{sec:discussion}.
Appendix~\ref{app:audit_checklist} provides a compact audit checklist of the logical dependencies.

\subsection{Contributions and outline}
\label{sec:intro_contributions}

The paper makes four contributions.
\begin{itemize}
\item \textbf{Axiomatization of Wish.} We define Wish as protocol-stable period data and fix its protocol invariance and finite-resource realizability conditions.
\item \textbf{Auditable finite-resource certificates.} We present the on-protocol error decomposition into sampling discrepancy and regularization/truncation error and define a computable stability functional.
\item \textbf{Teleological variational dynamics on protocol space.} We introduce a certified discrepancy potential and derive a damped inertial gradient dynamics whose Lyapunov function encodes ``minimal computational discrepancy'' under a complexity penalty.
\item \textbf{Reproducible numerics.} We provide pure-Python scripts demonstrating scan realizations of $\log 2$, $\pi$, $\zeta(2)$, and $\zeta(3)$, a truncation error decomposition, a bounded-complexity search yielding a sharp uniqueness gap for a fine-structure baseline, and a toy simulation of the derived dynamics.
\end{itemize}

Section~\ref{sec:preliminaries} introduces protocols, readouts, period data, and the scan--period interface.
Section~\ref{sec:wish} defines Wish and states the realization theorem.
Sections~\ref{sec:error_certificates}--\ref{sec:selection_principle} present auditable error certificates and a falsifiable selection functional.
Section~\ref{sec:variational_dynamics} contains the main new closed contribution: variational dynamics on protocol space.
Sections~\ref{sec:constant_search}--\ref{sec:reproducible_experiments} provide a constant-rigidity toy model and reproducible experiments.
We conclude with a discussion of interpretation-layer mappings and falsifiable predictions in Section~\ref{sec:discussion}.


